\documentclass[11pt]{article}

\usepackage[margin=1in]{geometry}
\usepackage{amsmath,amssymb,amsthm}
\usepackage{booktabs}
\usepackage{enumitem}
\usepackage[colorlinks=true,linkcolor=blue,citecolor=blue,urlcolor=blue]{hyperref}

\newtheorem{theorem}{Theorem}[section]
\newtheorem{lemma}[theorem]{Lemma}
\newtheorem{proposition}[theorem]{Proposition}
\newtheorem{corollary}[theorem]{Corollary}
\newtheorem{claim}[theorem]{Claim}
\theoremstyle{definition}
\newtheorem{definition}[theorem]{Definition}
\newtheorem{principle}[theorem]{Principle}
\theoremstyle{remark}
\newtheorem{remark}[theorem]{Remark}

% Anchor macro: every structural claim in this essay carries one of these,
% pointing into the companion corpus.
\newcommand{\anchor}[1]{\textup{[#1]}}

\newcommand{\Z}{\mathbb{Z}}
\newcommand{\N}{\mathbb{N}}
\newcommand{\Ztwo}{\mathbb{Z}_2}
\newcommand{\vtwo}{v_2}
\newcommand{\vthree}{v_3}
\newcommand{\Info}{\mathrm{I}}
\newcommand{\Ent}{\mathrm{H}}
\newcommand{\CV}{\mathrm{CV}}

\title{The $3x+1$ problem as interface physics:\\
two substrates, quantized resonance, and the limits of proof}

\author{Martien de Jong \and Jengo (AI collaboration)}

\date{July 2026}

\begin{document}

\maketitle

\begin{abstract}
This is a discussion paper, but of an unusual kind: every structural claim in it
is backed by a theorem, an exact identity, or an exhaustive computation from a
companion corpus on the $3x+1$ map \cite{Note2026, Forgetful2026, CstComma2026,
Tempering2026, Drift2026, GammaOne2026}. We propose reading the map as the
\emph{minimal coupling of two substrates}: a $2$-adic substrate in which halving
is clean, and a $3$-adic substrate in which the odd step is a pure digit append
--- the unique carry-free member of the $3n+d$ family. Each substrate alone is
completely integrable; all difficulty is concentrated at the interface, whose
coupling constant is the irrational relation $2^{\alpha}=3$. From this reading
four theses follow, each anchored in proof rather than analogy. (1) The map's
celebrated pseudo-randomness is a \emph{theorem}, not an observation: refills
are exactly fair, the macro-step law factorizes exactly, storage is exactly
zero, and cross-base mutual information vanishes --- randomness is what one
substrate's observer sees of clockwork coherent in the other, and noise is a
relation between incommensurable receivers, not a property of the signal.
(2) Cross-substrate coherence is not impossible but \emph{quantized}: cycles
are standing waves living only on the ladder of continued-fraction convergents
of $\log_2 3$ (Arnold tongues of the $2{:}3$ coupling), a correspondence made
flesh by the verified cycle zoo of the $3n+d$ family ($3n+13$, sitting on the
Pythagorean limma $256/243$, has exactly seven free $5$-step cycles).
(3) Wolfram's irreducibility thesis must be refined in both directions: the
ensemble level is \emph{more} reducible than his picture suggests (exactly
solvable), while the pointwise level resists all reduction \emph{provably};
moreover the zero-storage results weigh against the standard (Conway-style)
mechanism for undecidability of this instance. (4) The two failure modes of the
conjecture obey opposite ledgers --- cycles are localized and spectral,
divergence is delocalized and diffuse --- and this asymmetry is mirrored in
logic: membership in the tree of $1$ has short certificates of length
$O(\log n)$, while non-membership is not even semi-decidable by computation.
The hypothetical ``infinite winner'' demanded by compactness intuition does
exist --- but in $\Ztwo$, not in $\Z$: the conjecture is exactly the statement
that the limit always escapes the integers. We close with an honest ledger of
what is proved, what is structural template, and the three walls that remain.
\end{abstract}

\section{The two-substrate reading}
\label{sec:substrates}

\subsection{A map with two homes}

The Syracuse form of the Collatz map sends an odd integer $n$ to
$T(n) = (3n+1)/2^{\vtwo(3n+1)}$. Sixty years of work have treated this as one
dynamical system on one space. The reading we develop here --- and which the
companion corpus turns from metaphor into mathematics --- is that it is better
seen as \emph{two} exactly solvable systems and one hard \emph{interface}.

Write an integer simultaneously in base $2$ and in base $3$. The two halves of
the map then live in different worlds:

\begin{itemize}[leftmargin=2em]
\item \textbf{The $2$-adic substrate.} Halving is a pure shift in binary: the
binary digit sum $s_2$ is invariant, and the branching decision (how many
halvings $w$ follow an odd step) is read off the low binary bits. The map's
entire control flow --- the parity word --- is $2$-adic. By the Terras
correspondence, depth-$r$ index patterns occupy exactly one residue class
modulo $2^{\text{bits used}}$, so pattern densities are \emph{exact}, not
asymptotic \anchor{Note, Thm.~10}, \cite{Terras1976}.
\item \textbf{The $3$-adic substrate.} The odd step is invisible complexity in
base $2$ (a multiplication with carries) but trivial in base $3$: $3r+1$
\emph{appends the trit} $1$, with zero propagation into the digits of $r$, so
the ternary digit sum obeys $s_3 \mapsto s_3+1$ exactly \anchor{Note,
Thm.~25}. Halving in base $3$ is likewise an exact automaton: a top-down
parity sweep with a six-entry shed table \anchor{Note, Thm.~36}, and even the
branch condition closes inside base $3$, since $x \bmod 2^j$ is an explicit
periodic-weight trit sum \anchor{Note, Lem.~37}.
\end{itemize}

So each substrate, taken alone, is clean: the map is an odometer in one world
and a clockwork in the other. Where, then, does sixty years of difficulty
live?

\subsection{Each substrate alone is exactly solvable --- and therefore empty}

The corpus contains two no-go results, one per substrate, and together they
answer the question with unusual precision.

On the $2$-adic side: \emph{no function of the residue $n \bmod 2^j$ can be a
strict Lyapunov certificate at any depth $j$} \anchor{Note, Thm.~11}. The
proof is the sharpness of the shadowing bound: the residue class of $-5$ mod
$2^j$ admits, at every depth, a consistent macro-step continuation returning
to the class of $-5$ with value multiplied by $9/8 > 1$. The negative cycles
of the map are proof obstructions for the whole class of congruence-state
methods; any certificate must use archimedean information.

On the $3$-adic side: after an entire campaign that produced the exact
division automaton and a family of exact conservation laws --- the top tax
\anchor{Note, Thm.~40}, the zeros ledger \anchor{Note, Thm.~41}, run pairing
\anchor{Note, Thm.~44}, the slot lifecycle \anchor{Note, Thm.~43} --- the
sector was proved \emph{completely integrable}: every static ledger and every
digit-sum currency in the append/sweep automaton closes tautologically by
conservation, and no Lyapunov function exists in this sector either
\anchor{Note, Prop.~45}.

This pair of results deserves to be stated as the structural heart of the
two-substrate reading:

\begin{principle}[Integrable substrates, hard interface]
\label{prin:interface}
Each substrate of the $3x+1$ map is completely integrable, and precisely for
that reason carries no proof: its conservation laws close on themselves. All
content of the conjecture lives at the interface between the substrates,
whose coupling constant is the irrational relation $2^{\alpha} = 3$
(i.e.\ $\alpha = \log_2 3 \notin \mathbb{Q}$).
\end{principle}

The phrase ``coupling constant'' is not decorative. The identity $2^\alpha=3$
is literally what pins the interface to criticality: at the edge of the
Krasikov--Lagarias system the average row mass is \emph{exactly} $1$ because
$3/4 + 3\cdot 2^{\alpha-2} = 3$ \anchor{Note, Lem.~24}, and the offset-mass
flows up and down the $3$-adic scale ladder balance exactly, $3/16 = 3/16$
\anchor{Note, Prop.~32(i)}. The defining equation of the problem is precisely
the statement that the system sits on the conservative line: neither substrate
wins for free.

\subsection{Minimality: $3n+1$ is the weakest coupling that still couples}

Within the family $3n+d$ ($d$ odd), the choice $d=1$ is distinguished twice
over, and both distinctions are theorems.

First, the carry characterization: $3n+1$ is the \emph{unique} base-$3$-local
member of the family. For every other offset, $3r+d$ carries or borrows into
the digits of $r$ with geometrically distributed depth (for $3r-1$, the borrow
depth law is exactly $P(\text{depth}=j) = 2/3^j$) \anchor{Note, Thm.~25}. Only
$d=1$ touches the ternary clockwork with a bare append --- the minimal
possible contact.

Second, the pair law: consecutive segment endpoints within a family satisfy
$x' = 3x+2$, hence $x'/2 = 3(x/2)+1$; the guaranteed merge structure of the
Collatz graph exists for the multiplier $3$ and for \emph{no other} map $cn+1$
\anchor{Note, Thm.~2}.

The convergent map is thus the minimally coupled and maximally local member of
its family. This inverts the usual intuition that $3n+1$ is a generic
representative made hard by accident: it is the \emph{extremal} member, the
weakest coupling that still connects the two substrates --- and, empirically
and conjecturally, the one that converges. The reading suggested by the whole
corpus is that these two facts are the same fact.

\section{Randomness from determinism as a theorem}
\label{sec:randomness}

\subsection{Three grades of deterministic pseudo-randomness}

Wolfram's \emph{A New Kind of Science} \cite{Wolfram2002} made Rule 30 the
emblem of randomness generated by simple deterministic rules. But the epistemic
status of that randomness is empirical: the center column of Rule 30 passes
statistical tests, and no proof of any strong randomness property is known.
Rule 110 sits differently: its interesting property --- universality --- is
proved \cite{Cook2004}, and its randomness is incidental. The $3x+1$ map, we
claim, occupies a third grade, strictly above both for this purpose:
\emph{its randomness is a theorem}. Specifically:

\begin{itemize}[leftmargin=2em]
\item \textbf{Exact fairness.} The set of odd $n$ whose first $r$ successive
ladder refills (trailing-ones depths after each full burn) are all $\geq j$
has density \emph{exactly} $(2^{-(j-1)})^r$: successive refills are exactly
independent fair geometric draws at density level \anchor{Note, Thm.~20}.
\item \textbf{Exact factorization.} Exhaustively over odd $n < 2^{22}$, the
within-step joint law of (trailing ones $k$, post-burn halvings $w$)
factorizes exactly: $\Info(k;w) = 0.000000$ bits and $P(w \mid k) = 2^{-w}$
for every $k$, so $P(k,w) = 2^{-k}2^{-w}$ \anchor{Note, Thm.~27}.
\item \textbf{Zero storage.} The mutual information between refill depths
along macro-orbits, $\Info(k_0; k_d)$ for $d=1,2,3$, is $0$ to within
$3\times 10^{-6}$ bits against an entropy of $1.9375$ bits: the macro-automaton
has zero channel capacity at density level \anchor{Note, Thm.~22}.
\item \textbf{Cross-base blindness.} The apparent coupling
$\Info(n \bmod 9;\, w_{\mathrm{next}})$ vanishes as the $2$-adic modulus
grows: it was finite-modulus leakage, and the true cross-base mutual
information is zero \anchor{Note, Thm.~27}.
\end{itemize}

These are counted quantities, not estimated ones. The density-level stochastic
model of the map is \emph{exactly} independent fair geometric dice with no
measurable correction at any level; every hiding place for density-level
structure is closed \anchor{Note, Thm.~27; Forgetful}.

\subsection{The mechanism: noise as inter-receiver incommensurability}

What makes this more than a curiosity is that the corpus also identifies the
\emph{mechanism}, and the mechanism is the interface of
Section~\ref{sec:substrates}. The ``randomness'' seen in the $2$-adic frame is
the $2$-adic reading of structure that is perfectly coherent in the $3$-adic
frame: an append-only clockwork \anchor{Note, Thm.~25}, driven at the top by a
Sturmian rotation governed by the irrationality of $\log_2 3$ \anchor{Note,
Prop.~38}. Below the interface sits a rational odometer (perfectly periodic
bit-clocks); above it, an irrational rotation; their product, read in either
single-base frame, is exact fair dice. The incommensurability of the coupling
constant $2^\alpha = 3$ \emph{is} the noise source.

This suggests a definition that we state as a proposal, because its
mathematical content here is proved while its generality is not:

\begin{principle}[Noise as a relation]
\label{prin:noise}
Noise is not a property of a signal; it is a relation between receivers.
Determinism in substrate $A$ appears as randomness in receiver $B$ exactly
when their bases are incommensurable. In the $3x+1$ instance this is
quantitative: the fairness is exact at ensemble level \anchor{Note,
Thms.~20, 27}, and the anti-storage theorem \anchor{Note, Thm.~22} shows no
bounded single-base receiver can exploit the determinism underneath ---
there is no channel --- while the full two-base observer sees pure clockwork.
\end{principle}

Under this reading the conjecture itself acquires a striking form: even the
orbit's \emph{own thread} cannot exploit its determinism against the
interface. Convergence is self-decoherence.

\section{Quantized resonance: cycles on the comma ladder}
\label{sec:resonance}

\subsection{Not ``no coherence'' but ``spectral coherence''}

Section~\ref{sec:randomness} might suggest that cross-substrate coherence is
simply impossible: exact fair dice, zero storage, zero mutual information.
But that would overstate what is proved, and the overstatement would miss the
most beautiful structure in the problem. What the theorems exclude is
\emph{generic}, broadband coherence: a typical signal cannot hold both
receivers. Coherent objects nonetheless exist --- cycles. A cycle is a
standing wave at the interface: periodic, hence coherent in the $2$-adic frame
(its parity word repeats) and in the $3$-adic frame (its trit-stream
statistics repeat), forever.

Where can such standing waves live? The cycle equation for a shape with $k$
odd steps and $s$ total halvings reads
\[
n \,(2^{s} - 3^{k}) \;=\; B(\text{shape}), \qquad B > 0
\text{ a sum of positive terms,}
\]
so an integer cycle forces $2^s - 3^k$ to be small and positive relative to
$B$: the ratio $s/k$ must be a near-best rational approximation of
$\alpha = \log_2 3$ \anchor{Note, Thm.~5}. The permitted sites are the
continued-fraction convergents of $\log_2 3$ --- in musical language, the
ladder of \emph{commas}: $2^2/3^1$ (the touch $4-3=1$), $2^5/3^3 = 32/27$,
$2^8/3^5 = 256/243$ (the Pythagorean limma), $2^{19}/3^{12}$ (the comma
proper), and so on. Cycles are not forbidden; they are \emph{quantized}. The
dynamical-systems name for this phenomenon is mode-locking: the comma ladder
is the Arnold-tongue spectrum \cite{Arnold1965} of the $2{:}3$ coupling.

\subsection{The comma--cycle correspondence, made flesh}

If this spectral picture is right, it should be \emph{constructive}: sitting a
map exactly on a rung of the ladder should produce cycles in abundance. It
does, and this is a theorem with a verified zoo attached. For each near-touch
pair $(3^k, 2^s)$ with gap $d = 2^s - 3^k > 0$, the map $n \mapsto (3n+d)/2^w$
has \emph{free} cycles with exactly $k$ odd steps and $s$ halvings: the cycle
equation $n(2^s-3^k) = dB$ cancels to $n = B$, so every composition shape
yields a cycle, subject only to parity --- $\binom{s-1}{k-1}$ shapes per rung
\anchor{Note, Prop.~48}. Verified instances:

\begin{center}
\begin{tabular}{llll}
\toprule
map & rung of the ladder & gap $d$ & free cycles (odd starters) \\
\midrule
$3n+7$ & $2^4/3^2 = 16/9$ & $7$ & $5,\ 7,\ 11$ \\
$3n+5$ & $2^5/3^3 = 32/27$ & $5$ & $19,\ 23,\ 29,\ 31$ \\
$3n+13$ & $2^8/3^5 = 256/243$ (limma) & $13$ &
$211,\ 227,\ 251,\ 259,\ 283,\ 287,\ 319$ \\
\bottomrule
\end{tabular}
\end{center}

\noindent The map $3n+13$, sitting on the limma, has exactly \emph{seven}
distinct free $5$-step cycles; explicitly,
$211 \to 323 \to 491 \to 743 \to 1121 \to 211$. The notorious cycle zoo of the
$3n+d$ family is not a zoo at all: it is the comma ladder made flesh, each
convergent of $\log_2 3$ endowing its gap value $d$ with a full resonance
family.

\subsection{Why $3n+1$ is poor: the free resonances are exhausted}

Conversely, the cycle \emph{scarcity} of $3n+1$ is the statement that $d=1$
sits only on the exact touches: integrality in the cycle equation is automatic
iff $|2^{s} - 3^{k}| = 1$. By Mih\u{a}ilescu's theorem (Catalan's conjecture)
\cite{Mihailescu2004} the difference-$1$ pairs are exhausted by
$(k,s) = (1,2)$ and $(2,3)$: exactly the trivial cycle $\{1,4,2\}$ and the
$-5$ cycle. The two free cycles arithmetic can give, it has given
\anchor{Note, Thm.~5, Cor.}.

For the cases actually needed, full Catalan is a luxury: elementary
congruences suffice, and the observation is old. Modulo $8$, $3^m + 1$ is
$2$ or $4$, never $0$, so $2^n = 3^m + 1$ has no solution with $n \geq 3$
--- one line kills all infinitely many candidates at once \anchor{Note,
\S12}. And $3^m = 2^n + 1$ forces $m$ even, whence
$(3^{m/2}-1)(3^{m/2}+1) = 2^n$ gives $9 = 8+1$ only. This last case --- that
$(8,9)$ is the final pair of consecutive $3$-smooth numbers --- was proved by
Levi ben Gershon (Gersonides) in 1343, in \emph{De numeris harmonicis},
answering a question of the composer Philippe de Vitry about which pairs of
harmonic numbers can differ by one \cite{Gersonides1343}. The exhaustion of
the free resonances of the $2{:}3$ interface is, fittingly, a result born in
music theory --- the same ladder that governs why a circle of twelve fifths
misses seven octaves by a comma.

The ground state $\{1,4,2\}$ is itself a resonance, living on the touch
$4 - 3 = 1$. The conjecture is therefore \emph{not} ``there is no resonance''
--- that is false, the trivial cycle is one --- but:

\begin{principle}[The conjecture as a ground-state statement]
The ground resonance is globally attracting: every orbit decoheres into the
one permitted mode. Between incommensurable substrates, coherence exists only
on the discrete comma spectrum of the coupling constant; broadband coherence
is impossible; and released systems drain into the lowest available
resonance.
\end{principle}

\subsection{The same ladder governs the stopping-time structure}

That the comma ladder is the problem's true skeleton --- not merely its cycle
bookkeeping --- is confirmed from an independent direction. Terras'
coefficient-stopping-time conjecture (that $\tau(n) = \sigma(n)$ for $n>1$)
was verified exhaustively for all stopping classes with $\tau \leq 24$
($81{,}119$ classes) and all odd $1 < n \leq 10^7$; the smallest safety margin
observed is $2.02$, and the extremal classes concentrate at
$(u,j) = (5,8)$ --- the limma again \anchor{Note, Props.~28, 29}. A CST
violation in class $(u,j)$ requires $n \leq b/(2^j - 3^u)$, which explodes
exactly at the convergents of $\log_2 3$; with the five-line bound
$b/2^j < u/3$ the reduction becomes unconditional and extends CST to
$\tau \leq 4700$ \anchor{CstComma, Cor.~6.2}. Even the growth of the extremal
numerators is comma-governed: $B(t) \sim 1 + t/8$, one unit per repetition of
the limma word $(5,8)$ \anchor{Note, Lem.~30}. Cycle exclusion and CST are
one wall seen from two sides: effective rational approximation of $\log_2 3$
--- Baker territory \cite{Baker1966, Rhin1987}.

\section{Wolfram's irreducibility thesis, refined}
\label{sec:wolfram}

Wolfram presents the Collatz problem as a poster child of computational
irreducibility, with the suggestion --- backed by Conway's universality
theorem for the generalized class \cite{Conway1972} --- that the instance may
be undecidable \cite{Wolfram2002}. The corpus forces a refinement of this
picture in \emph{both} directions at once, and the refinement is more
interesting than either verdict alone.

\subsection{The ensemble level is more reducible than Wolfram's picture}

Irreducibility, in Wolfram's usage, admits ``pockets of reducibility.'' What
the corpus shows is stronger: at ensemble (density) level the $3x+1$ map is
not a pocket but \emph{exactly solvable}. The stationary roulette measure has
closed form (mod $9$: $\pi(1,2,4,5,7,8) = (8,16,11,4,2,22)/63$); refills are
exactly fair \anchor{Note, Thm.~20}; the macro-step law factorizes exactly
\anchor{Note, Thm.~27}; storage is exactly zero \anchor{Note, Thm.~22}; the
linearized Krasikov--Lagarias edge operator has spectrum
$\{1\} \cup \{|\lambda| = 1/4\}$ exactly \anchor{Note, Thm.~21}; and the
Perron eigenvectors of the certificate system are, to $R^2 > 0.99$, pure
powers of the roulette measure \anchor{Tempering}. In the language of
coarse-graining: lumpability of the macro-chain is exact --- a theorem here,
not an approximation. Tao's celebrated ``almost all orbits attain almost
bounded values'' \cite{Tao2019} is, in this reading, exactly what
coarse-graining success at ensemble level looks like from the analytic side.

\subsection{The pointwise level resists all reduction, provably}

At the level of the individual orbit the resistance is not an impression but
a set of impossibility results: no residue-based Lyapunov certificate at any
$2$-adic depth \anchor{Note, Thm.~11}; complete integrability (hence
certificate-emptiness) of the entire digit-sum sector \anchor{Note,
Prop.~45}; mirror-blindness --- the Krasikov--Lagarias growth exponent is
identical for $3n-1$ and $3n+1$, so density methods cannot even distinguish
the convergent map from a map with three known nontrivial cycles
\anchor{Forgetful}; and the no-go on coarse resampling --- the cascade
attenuation $\kappa$ is not a functional of any coarse histogram; only exact
structure reproduces it \anchor{Note, Campaign~XVIII}. The
reducible/irreducible boundary is thus not fuzzy but \emph{located}: it is
the ensemble/pointwise split, and it is proved from both sides.

\subsection{Against the hidden computer}

Wolfram's suspicion of undecidability inherits its plausibility from Conway:
generalized Collatz maps can simulate arbitrary register machines, encoding
registers in the exponents of primes, and the generalized problem is
$\Pi^0_2$-complete \cite{Conway1972, KurtzSimon2007}. But every known
construction of this type works by \emph{storing information}. The specific
$3n+1$ instance exhibits the opposite signature on every test: zero channel
capacity at density level \anchor{Note, Thm.~22}, exact memorylessness
\anchor{Note, Thms.~20, 27}, and the carry characterization --- the unique
append-only, maximally forgetful member of its family \anchor{Note, Thm.~25}.
This is the signature of an odometer, not of a Turing machine
\anchor{Forgetful}. Undecidability of the instance remains possible; what the
evidence weighs against is the \emph{standard mechanism} for it. If the
$3x+1$ problem is independent of strong theories, it will not be because a
computer hides in it, but for a reason not yet named --- which would itself
be a discovery about the reach of the Conway paradigm.

\section{The two requirement ledgers and the semi-decidability asymmetry}
\label{sec:ledgers}

The conjecture has exactly two failure modes: a nontrivial cycle, or a
divergent orbit. The corpus compiles, from proved results only, what each
would have to satisfy --- and the two ledgers turn out to have \emph{opposite
shapes}.

\subsection{The cycle ledger: localized and spectral}

A hypothetical cycle with $k$ odd steps and $s$ halvings must, simultaneously
\anchor{Note, \S11 of the synthesis; Thms.~5, 9, 17}:
\begin{enumerate}[leftmargin=2em, itemsep=1pt]
\item live on the comma ladder: $n = B/(2^s-3^k) \in \Z$ forces
$s/k \approx \log_2 3$ (near-touch), so its possible homes are spectrally
enumerated;
\item survive the exhaustion of the free homes: gap-$1$ sites carry only the
trivial cycle (elementary today; Gersonides 1343 for the $9/8$ case;
Mih\u{a}ilescu in general), so it needs the divisibility miracle
$(2^s-3^k) \mid B(\text{shape})$ against superpolynomially growing gaps;
\item have a net-falling shape, by the Sign Theorem: every net-climbing shape
realizes at a \emph{negative} $2$-adic point \anchor{Note, Thm.~5};
\item balance exact periodic ledgers: per period the digit-sum shed equals
$u$ exactly and the top tax equals $j$, so its lower sweep must gain exactly
$j-u$ per period \anchor{Note, Thm.~40};
\item be large: no integer cycles with $\leq 14$ odd steps beyond the known
four $\{1, -1, -5, -17\}$ \anchor{Note, Thm.~17, extended}; at least
$1.375 \times 10^{11}$ odd terms and $\geq 92$ circuits by Hercher's bound
\cite{Hercher2023}; all elements above $2^{71}$ by exhaustive verification
\cite{Barina2021}.
\end{enumerate}
This ledger is \emph{localized}: the candidate sites are known, discrete,
and individually checkable; finitely many checks per site plus one effective
gap bound (Baker-type) would finish cycles off entirely. A cycle, if it
exists, is a standing wave on a named Arnold tongue --- we know exactly where
it could live.

\subsection{The divergence ledger: delocalized and diffuse}

A hypothetical divergent orbit must \anchor{Note, \S11; Thms.~9, 9$'$, 10, 20,
40; Prop.~13}:
\begin{enumerate}[leftmargin=2em, itemsep=1pt]
\item hold $\bar{w} < \log_2 3$ forever --- a sustained excess of $w=1$ steps
against dice proved exactly fair, an event of measure zero in every window,
required pointwise \anchor{Note, Thm.~20};
\item have an aperiodic parity word (else it is a cycle), and moreover escape
\emph{every} periodic and eventually-periodic index pattern within
logarithmically many digits: shadowing any non-realized rational phase is
throttled at $r \leq \log_2(n+|\rho|)/d$ periods, unconditionally
\anchor{Note, Thms.~9, 9$'$} --- divergence requires unbounded aperiodic
complexity, quantitatively;
\item pay for every monotone climb at the exact exponential price
$\mu^r$ with $\mu = 0.28627450\ldots$ \anchor{Note, Thm.~10};
\item pay the top tax --- an exact $+1$ shed at the top of every halving,
pointwise, every orbit \anchor{Note, Thm.~40} --- while its length grows;
\item and yet need \emph{no arithmetic coincidence at all}: no divisibility
miracle, no special site. Its requirement is diffuse --- it must win
everywhere, forever.
\end{enumerate}
Divergent orbits are excluded to density $x^{0.9146}$ (certified at $3$-adic
depth $k=20$ \anchor{Tempering}) and below $2^{71}$ \cite{Barina2021}; what
survives is a measure-zero possibility with no location. The asymmetry is
stark: \emph{cycles are localized and spectral; divergence is delocalized and
diffuse}. The wall protecting cycles is transcendence theory; the wall
protecting divergence is the pointwise/measure-zero gap itself.

\subsection{The logical mirror: membership confesses, non-membership never}

This dynamical asymmetry has an exact logical shadow. Membership in the tree
of $1$ is semi-decidable with \emph{short} certificates: if $n$ reaches $1$,
its own orbit is the proof, and its length is at most $c^*(n)\log_2 n$ with
the laggard-edge constant $c^* \approx 18$ at $2^{24}$, growing slowly toward
$\approx 29$ --- a polynomial-size certificate \anchor{Note, \S13}. Every
number ever named has been proven in; there are no known hard numbers, only
numbers nobody has run. If the conjecture is true, \emph{every} integer has a
short membership proof.

Non-membership is worse than hard: it is not even semi-decidable by running,
because a divergent orbit never confesses. No individual integer can ever be
proven \emph{out} by computation alone. Uniformly, the residue classes
lacking a depth-$k$ proof are exactly the survivor classes mod $2^k$ ---
counting $\sim 1.834^k$, density $0.917^k \to 0$, ones-rich (the kin of $63
\bmod 64$, the shadow set of the rational phases $-1, -5, -17$, mean
ones-fraction $0.653$) \anchor{Note, \S13}. At proof budget $t$ the unproven
region is exactly the wavefront $\log_2 n > t/c^*$: the tree is provable
number by number at logarithmic cost each; what cannot be bought at any
finite budget is the single sentence \emph{and this works for all of them}.
This is the $\Pi^0_2$ shape of the problem \cite{KurtzSimon2007} made
concrete: a $\Sigma_1$ membership predicate with short witnesses, universally
quantified.

\section{Where the infinite winner lives}
\label{sec:winner}

There is a persistent intuition --- every student of the problem has felt it
--- that runs: \emph{there are infinitely many integers; the conditions for
escape are just statistics; surely some number satisfies all of them}. The
corpus lets us answer this intuition precisely, and the answer has three
layers, of which the third is the deepest fact in this essay.

\textbf{(a) For cycles, the conditions are equations, not lottery tickets.}
Infinitely many candidates are powerless against a modular obstruction: there
are infinitely many integers, and none of them satisfies $2^n = 3^m+1$ beyond
$4 = 3+1$ --- one congruence mod $8$ kills them all at once \anchor{Note,
\S12}. For the remaining candidate sites, the expected number of divisibility
miracles is $\sum_{\text{sites}} O(1/\text{gap})$, which \emph{converges}:
infinity supplies candidates logarithmically, the conditions cut
exponentially.

\textbf{(b) For divergence, the heuristic argues for the conjecture.} Per
candidate $n$, the model probability of beating the exactly fair coin forever
is the limit of $2^{-cT}$, which is $0$; countable additivity over infinitely
many candidates gives expected winners $= 0$. Infinity times zero is zero
here, not a loophole.

\textbf{(c) And yet the winner exists --- one substrate up.} At every finite
depth $T$, integers satisfying the first $T$ escape conditions \emph{do}
exist: the survivor classes are never empty, numbering $S(T) \sim 1.834^T$
\anchor{Note, \S13}. But they are \emph{different} integers at each depth,
running off to infinity, and the sequence of satisfiers converges --- in
$\Ztwo$, not in $\Z$. The pure-climb example is exact: the satisfier of the
$k$-fold climb condition is $2^k - 1$, and as $k \to \infty$,
\[
2^k - 1 \;\longrightarrow\; -1 \;=\; \ldots 111_2 \ \in \Ztwo,
\]
a genuine fixed point of pure climbing --- and not a positive integer. In the
$2$-adic completion divergent orbits are not rare but \emph{abundant}: almost
every $2$-adic point diverges. The compactness winner your intuition demands
exists as a pseudo-orbit: real in one receiver, noise in the other, finite in
one substrate only.

This is where the two-substrate reading closes its loop. The integers are
precisely the objects that are \emph{finite in both substrates
simultaneously} --- finite binary top and finite ternary top; $2$-adic
pseudo-orbits are finite in one receiver only, and the archimedean property
is exactly dual finiteness \anchor{Note, \S3 of the synthesis}. Descent is
\emph{powered} by the finite top: the digit-sum drainage that pays the
conservation ledger is produced almost entirely at the top boundary
($+0.485$ per trit in the top decile, against $\approx 0$ in the bulk), the
one feature integers have and pseudo-orbits lack \anchor{Note, Props.~38, 39,
Thm.~40}. Hence:

\begin{principle}[The conjecture as an escape statement]
The survivor construction always produces its winner --- in $\Ztwo$. The
Collatz conjecture is exactly the statement that the limit always escapes the
integers: no dually-finite point realizes the infinite climbing stream. The
infinite winner demanded by compactness intuition lives one substrate up,
outside the countable thread.
\end{principle}

\section{Honest closing: proved, template, and the three walls}
\label{sec:closing}

A synthesis essay owes its reader an exact accounting of epistemic status.
We close with three lists.

\subsection{What is proved or exhaustively computed}

Every load-bearing claim above is anchored: the carry characterization and
uniqueness of $3n+1$ \anchor{Note, Thm.~25}; the pair-law uniqueness of the
multiplier $3$ \anchor{Note, Thm.~2}; the two substrate no-gos \anchor{Note,
Thm.~11, Prop.~45}; exact fairness, factorization, and zero storage
\anchor{Note, Thms.~20, 22, 27}; the conservative-line identities pinned by
$2^\alpha = 3$ \anchor{Note, Lem.~24, Prop.~32}; the Sign Theorem and the
Catalan corollary \anchor{Note, Thm.~5}; the cycle census through period $14$
\anchor{Note, Thm.~17}; the shadowing bounds \anchor{Note, Thms.~9, 9$'$};
the comma--cycle correspondence with its verified zoo \anchor{Note,
Prop.~48}; the CST reduction to the convergents of $\log_2 3$, unconditional
to $\tau \leq 4700$ \anchor{CstComma}; the certified density exponent
$\gamma = 0.9146$ \anchor{Tempering}; and the exact edge rate
$d\gamma/dq = 1/\ln(4/3)$ \anchor{Note, Thm.~19}. Where the corpus retracted
--- the per-position frontier ledger, the flagged $\bar{w} \leq 1.226$ bound
--- the retractions are recorded in the same documents \anchor{Note,
corr.~to Thm.~42, Prop.~45}, and nothing in this essay leans on them.

\subsection{What is structural template, not theorem}

The physics vocabulary --- receivers, coherence, decoherence, resonance,
mode-locking --- is proved \emph{about the $2{:}3$ arithmetic interface} and
about nothing else. The proposal that noise in general is inter-receiver
incommensurability (Principle~\ref{prin:noise}), and any transfer of the
mode-locking template to interfaces outside arithmetic, is a structural
analogy whose merit is that it is falsifiable in form: a claimed instance
must name its coupling constant, its convergent, and its tongue width. The
mathematics here cannot say whether such transfers are deep or merely
well-matched; we flag the boundary and stay on the proved side of it.

\subsection{The three walls}

What remains of the $3x+1$ problem, after the corpus, is exactly three
walls --- and naming them precisely is the essay's final claim.

\begin{enumerate}[leftmargin=2em]
\item \textbf{The transcendence wall} (cycles, and CST). Both cycle exclusion
and the coefficient-stopping-time conjecture reduce to effective lower bounds
on $|2^j - 3^u|$ at the convergents of $\log_2 3$ \anchor{Note, Prop.~28,
Thm.~29}: one wall, two classical sub-problems leaning on it. Current
effective irrationality measures \cite{Baker1966, Rhin1987} push the danger
zone out (first convergent escape at $(u,j) = (15601, 24727)$
\anchor{CstComma}) but do not close it.
\item \textbf{The pointwise wall} (divergence). Every density-level statement
is now exact, and exactly fair --- and therefore exactly powerless per orbit:
the pointwise version of Theorem 20's fairness \emph{is} the remaining
content of the divergence half \anchor{Note, Thm.~20}. The sharpest exact
formulation is the automaton form: no finite-top trit stream can sustain the
self-referential $(c,d)$-statistics that perpetual climb requires
\anchor{Note, Thms.~36, 40, Lem.~37} --- the archimedean bridge, still
unbuilt.
\item \textbf{The contraction wall} (density $\gamma \to 1$). The
Krasikov--Lagarias program has been reduced to the stability of a
one-dimensional fixed point: prove the cascade attenuation $\kappa$ stays
uniformly below $1$ (measured $0.839 \pm 0.002$, flat in depth) and
$\gamma \to 1$ follows through proved identities \anchor{Note, Thm.~19,
Props.~35, 47; Drift; GammaOne}. This is the smallest formulation the
density problem has ever had --- and even fully won, it buys $x^{1-\epsilon}$,
not the pointwise sentence.
\end{enumerate}

The three walls are not equally shaped: the first is localized (finitely many
sites plus one effective bound), the third is a fixed-point stability
statement with every coefficient measured, and the second is the conjecture's
irreducible core --- the place where, on the evidence assembled here, the
ensemble world ends and the individual thread begins. The two-substrate
reading does not breach that wall. Its contribution is to show exactly where
the wall stands, why everything on one side of it is clockwork, and why the
one sentence on the other side --- \emph{every dually-finite signal decoheres
to the ground state} --- is the whole of the $3x+1$ problem.

\subsection*{Acknowledgements}

This essay synthesizes results obtained over roughly 1200 rounds of
human--AI collaborative research; the division of labour is recorded in the
companion papers. The framing questions of Sections 2, 3, and 6 --- whether
the randomness has a mechanism, whether multi-layer resonance is impossible
or quantized, and where the infinite winner lives --- are due to the first
author.

\begin{thebibliography}{99}

\bibitem{Note2026}
M. de Jong and Jengo (AI),
\emph{Six results on the Collatz map in family/pair coordinates},
research note with continuing numbered corpus (Theorems 1--48), July 2026.
Cited as [Note].

\bibitem{Forgetful2026}
M. de Jong and Jengo (AI),
\emph{The $3x+1$ map is a maximally forgetful machine: exact fairness, zero
storage, and the carry characterization}, companion paper, July 2026.
Cited as [Forgetful].

\bibitem{CstComma2026}
M. de Jong and Jengo (AI),
\emph{The coefficient-stopping-time conjecture and the comma ladder:
reduction to effective rational approximation of $\log_2 3$}, companion
paper, July 2026. Cited as [CstComma].

\bibitem{Tempering2026}
M. de Jong and Jengo (AI),
\emph{Krasikov--Lagarias certificates for the $3x+1$ problem are tempered
roulette measures}, companion paper, July 2026. Cited as [Tempering].

\bibitem{Drift2026}
M. de Jong and Jengo (AI),
\emph{The drift of the Collatz cascade: the min-operator as a directional
low-pass filter}, companion paper, July 2026. Cited as [Drift].

\bibitem{GammaOne2026}
M. de Jong and Jengo (AI),
\emph{The $\gamma \to 1$ program: reduction of the density version of the
$3x+1$ problem to a one-dimensional fixed point}, companion paper, July 2026.
Cited as [GammaOne].

\bibitem{Arnold1965}
V. I. Arnold,
\emph{Small denominators. I. Mapping the circle onto itself},
Izv. Akad. Nauk SSSR Ser. Mat. 25 (1961), 21--86; Amer. Math. Soc. Transl.
(2) 46 (1965), 213--284.

\bibitem{Baker1966}
A. Baker,
\emph{Linear forms in the logarithms of algebraic numbers},
Mathematika 13 (1966), 204--216.

\bibitem{Barina2021}
D. Barina,
\emph{Convergence verification of the Collatz problem},
J. Supercomput. 77 (2021), 2681--2688; computation since extended to
$2^{71}$.

\bibitem{Conway1972}
J. H. Conway,
\emph{Unpredictable iterations},
Proc. 1972 Number Theory Conference, Univ. Colorado, Boulder (1972),
49--52.

\bibitem{Cook2004}
M. Cook,
\emph{Universality in elementary cellular automata},
Complex Systems 15 (2004), 1--40.

\bibitem{Gersonides1343}
Levi ben Gershon (Gersonides),
\emph{De numeris harmonicis} (1343); commissioned by Philippe de Vitry.
See J. Carlebach, \emph{Lewi ben Gerson als Mathematiker}, Berlin, 1910.

\bibitem{Hercher2023}
C. Hercher,
\emph{There are no Collatz $m$-cycles with $m \leq 91$},
J. Integer Seq. 26 (2023), Article 23.3.5.

\bibitem{KrasikovLagarias2003}
I. Krasikov and J. C. Lagarias,
\emph{Bounds for the $3x+1$ problem using difference inequalities},
Acta Arith. 109 (2003), 237--258.

\bibitem{KurtzSimon2007}
S. A. Kurtz and J. Simon,
\emph{The undecidability of the generalized Collatz problem},
Proc. TAMC 2007, Lecture Notes in Comput. Sci. 4484, Springer (2007),
542--553.

\bibitem{Lagarias1985}
J. C. Lagarias,
\emph{The $3x+1$ problem and its generalizations},
Amer. Math. Monthly 92 (1985), 3--23.

\bibitem{Lagarias2010}
J. C. Lagarias (ed.),
\emph{The Ultimate Challenge: The $3x+1$ Problem},
Amer. Math. Soc., Providence, RI, 2010.

\bibitem{Mihailescu2004}
P. Mih\u{a}ilescu,
\emph{Primary cyclotomic units and a proof of Catalan's conjecture},
J. Reine Angew. Math. 572 (2004), 167--195.

\bibitem{Rhin1987}
G. Rhin,
\emph{Approximants de Pad\'e et mesures effectives d'irrationalit\'e},
S\'emin. Th\'eor. Nombres, Paris 1985--86, Progr. Math. 71, Birkh\"auser
(1987), 155--164.

\bibitem{Tao2019}
T. Tao,
\emph{Almost all orbits of the Collatz map attain almost bounded values},
Forum Math. Pi 10 (2022), e12.

\bibitem{Terras1976}
R. Terras,
\emph{A stopping time problem on the positive integers},
Acta Arith. 30 (1976), 241--252.

\bibitem{Wolfram2002}
S. Wolfram,
\emph{A New Kind of Science},
Wolfram Media, Champaign, IL, 2002.

\end{thebibliography}

\end{document}
